% ============================================================================
% APPENDIX GA: LIT-AKERLOF - Identity Economics & Information Asymmetry
% ============================================================================
% Category: Literature Integration (LIT-)
% Language: English
% EBF Integration: Identity, asymmetric information, efficiency wages, animal spirits
% ============================================================================

\chapter*{Appendix GA: George Akerlof Research Integration}
\addcontentsline{toc}{chapter}{Appendix GA: LIT-AKERLOF - Identity Economics \& Information Asymmetry}
\label{app:akerlof}

% ============================================================================
% HEADER BLOCK
% ============================================================================
\begin{tcolorbox}[colback=blue!5!white, colframe=blue!75!black, title=Appendix Metadata]
\textbf{Code:} GA \\
\textbf{Category:} LIT- (Literature Integration) \\
\textbf{Title:} Identity Economics \& Information Asymmetry \\
\textbf{Version:} 1.0 \\
\textbf{Status:} Complete \\
\textbf{Primary Author Focus:} George Akerlof (Nobel Prize 2001) \\
\textbf{Papers Covered:} 30 \\
\textbf{Dependencies:} CORE-WHAT (C), CORE-AWARE (AU), DOMAIN-LABOR, CONTEXT-CULTURE \\
\textbf{EBF Connections:} Identity utility, information asymmetry, fairness, behavioral macro
\end{tcolorbox}

% ============================================================================
% CROSS-REFERENCE MAP
% ============================================================================
\begin{tcolorbox}[colback=green!5!white, colframe=green!75!black, title=Cross-Reference Map]
\textbf{Outgoing References:}
\begin{itemize}[nosep]
    \item CORE-WHAT (C): Identity as utility dimension
    \item CORE-AWARE (AU): Information asymmetry and awareness
    \item DOMAIN-LABOR: Efficiency wages and gift exchange
    \item DOMAIN-MACRO: Animal spirits and behavioral macro
    \item CONTEXT-CULTURE: Social norms and customs
    \item LIT-FEHR (K): Gift exchange experimental validation
\end{itemize}

\textbf{Incoming References:}
\begin{itemize}[nosep]
    \item Chapter 10 (Utility): Identity dimension formalization
    \item Chapter 9 (Context): Social norms as Ψ
    \item PREDICT-02: Identity-based predictions
\end{itemize}
\end{tcolorbox}

% ============================================================================
% ABSTRACT
% ============================================================================
\section*{Abstract}

George Akerlof's research has fundamentally shaped behavioral economics through four major contributions: \textbf{information asymmetry} (``The Market for Lemons''), \textbf{identity economics} (with Rachel Kranton), \textbf{efficiency wages and gift exchange}, and \textbf{behavioral macroeconomics} (``Animal Spirits'' with Robert Shiller). His 2001 Nobel Prize recognized work on asymmetric information. This appendix integrates Akerlof's 30 key papers into the EBF framework, with particular emphasis on:
\begin{enumerate}
    \item The \textbf{Lemons Problem} and adverse selection
    \item \textbf{Identity Economics}: How social categories affect utility
    \item \textbf{Gift Exchange}: Efficiency wages and fairness
    \item \textbf{Near-Rationality}: Small deviations with large consequences
    \item \textbf{Animal Spirits}: Confidence and behavioral macroeconomics
\end{enumerate}

% ============================================================================
% QUICK REFERENCE
% ============================================================================
\begin{tcolorbox}[colback=yellow!10!white, colframe=orange!75!black, title=Quick Reference: Key Akerlof Parameters]
\begin{tabular}{lll}
\textbf{Parameter} & \textbf{Value} & \textbf{Source} \\
\hline
Identity loss coefficient & $\delta = 0.2-0.5$ & Akerlof \& Kranton (2000) \\
Gift exchange reciprocity & $\gamma = 0.3-0.6$ & Akerlof (1982) \\
Efficiency wage premium & 10-20\% & Akerlof \& Yellen (1990) \\
Near-rationality threshold & $\epsilon < 0.01$ & Akerlof \& Yellen (1985) \\
Information asymmetry $\alpha$ & Context-dependent & Akerlof (1970) \\
\end{tabular}
\end{tcolorbox}

% ============================================================================
% FUNDAMENTAL QUESTION
% ============================================================================
% -----------------------------------------------------------------------------
% APPENDIX SCOPE BOX
% -----------------------------------------------------------------------------

\begin{tcolorbox}[
    colback=orange!5!white,
    colframe=orange!75!black,
    title={\textbf{Appendix Scope: Ziel / In-Scope / Out-of-Scope / Constraints}},
    fonttitle=\bfseries
]

\textbf{Ziel (Objective):}
\begin{itemize}[nosep]
    \item How do social identity, information asymmetries, and norms shape economic behavior and market outcomes beyond the standard rational actor model?
\end{itemize}

\textbf{In-Scope:}
\begin{itemize}[nosep]
    \item Fundamental Question
    \item Information Asymmetry: The Market for Lemons
    \item Identity Economics
    \item Gift Exchange and Efficiency Wages
    \item Near-Rationality and Behavioral Macro
\end{itemize}

\textbf{Out-of-Scope (delegiert an andere Appendices):}
\begin{itemize}[nosep]
    \item Glossary $\rightarrow$ Appendix G
\end{itemize}

\textbf{Constraints (Anwendungsgrenzen):}
\begin{itemize}[nosep]
    \item [TODO: Define when this appendix does NOT apply]
    \item [TODO: Define required prerequisites]
    \item [TODO: Define known limitations]
\end{itemize}

\textbf{Lieferobjekte (Deliverables):}
\begin{itemize}[nosep]
    \item 1 Table(s)
    \item 8 Equation(s)
    \item Worked Example(s)
\end{itemize}

\end{tcolorbox}


\section{Fundamental Question}

\begin{tcolorbox}[colback=red!5!white, colframe=red!75!black, title=Central Question]
\textbf{How do social identity, information asymmetries, and norms shape economic behavior and market outcomes beyond the standard rational actor model?}
\end{tcolorbox}

Akerlof's answer: Economic agents are embedded in \textbf{social categories} that prescribe appropriate behavior. When identity is threatened, agents sacrifice material payoffs. Information asymmetries can cause market failure. Small deviations from rationality can have large aggregate effects.

% ============================================================================
% THE MARKET FOR LEMONS
% ============================================================================
\section{Information Asymmetry: The Market for Lemons}

\subsection{The Core Problem}

Akerlof (1970) identified how asymmetric information causes market failure:

\begin{definition}[Adverse Selection]
When sellers have private information about quality:
\begin{enumerate}
    \item Buyers cannot distinguish high- from low-quality goods
    \item Buyers offer average price based on expected quality
    \item High-quality sellers exit (price too low)
    \item Average quality falls, price falls further
    \item Market may collapse entirely (``unraveling'')
\end{enumerate}
\end{definition}

\begin{axiom}[GA-1: Information Asymmetry]
\label{ax:ga1}
Market efficiency requires:
\begin{equation}
\alpha \equiv \frac{Information_{buyer}}{Information_{seller}} \geq \alpha^*
\end{equation}
When $\alpha < \alpha^*$, adverse selection dominates and markets fail.
\textbf{EBF Connection:} Extends CORE-AWARE with information completeness parameter.
\end{axiom}

\subsection{Solutions to the Lemons Problem}

\begin{itemize}
    \item \textbf{Signaling}: High-quality sellers invest in costly signals (warranties, brands)
    \item \textbf{Screening}: Buyers design mechanisms to separate types (insurance deductibles)
    \item \textbf{Institutions}: Third-party certification, reputation systems
    \item \textbf{Regulation}: Disclosure requirements, quality standards
\end{itemize}

% ============================================================================
% IDENTITY ECONOMICS
% ============================================================================
\section{Identity Economics}

\subsection{The Identity Framework}

Akerlof \& Kranton (2000, 2010) extended utility to include identity:

\begin{definition}[Identity Utility]
\begin{equation}
U_j = U_j(a_j, a_{-j}, I_j)
\end{equation}
where:
\begin{itemize}
    \item $a_j$ = own actions
    \item $a_{-j}$ = others' actions
    \item $I_j$ = identity utility derived from social categories
\end{itemize}
\end{definition}

\begin{axiom}[GA-2: Identity and Prescriptions]
\label{ax:ga2}
Identity utility depends on matching behavior to category prescriptions:
\begin{equation}
I_j = I_j(a_j, c_j, \mathbf{P})
\end{equation}
where $c_j$ is agent's social category and $\mathbf{P}$ is the vector of prescriptions for that category.
\textbf{EBF Connection:} Adds Identity (I) to CORE-WHAT utility dimensions (FEPSDIE).
\end{axiom}

\begin{axiom}[GA-3: Identity Loss]
\label{ax:ga3}
Deviating from category prescriptions causes identity loss:
\begin{equation}
\Delta I = -\delta \cdot \|a_j - P(c_j)\|^2
\end{equation}
where $\delta > 0$ is the identity sensitivity parameter.
\textbf{EBF Connection:} Quantifies cost of norm violation in CONTEXT-CULTURE.
\end{axiom}

\subsection{Applications of Identity Economics}

\begin{tcolorbox}[colback=blue!5!white, colframe=blue!75!black, title=Identity Economics Applications]
\begin{itemize}
    \item \textbf{Gender and Work}: ``Men don't do women's jobs'' reduces male nursing applicants
    \item \textbf{Education}: Student identity affects effort (``nerds'' vs. ``burnouts'')
    \item \textbf{Organizations}: Insider vs. outsider identity affects loyalty
    \item \textbf{Conflict}: In-group/out-group identity drives discrimination
\end{itemize}
\end{tcolorbox}

% ============================================================================
% GIFT EXCHANGE AND EFFICIENCY WAGES
% ============================================================================
\section{Gift Exchange and Efficiency Wages}

\subsection{The Gift Exchange Model}

Akerlof (1982) proposed that labor contracts involve reciprocal gift exchange:

\begin{definition}[Gift Exchange]
Firms pay above market-clearing wages as a ``gift''; workers reciprocate with above-minimum effort as a ``gift'' in return.
\end{definition}

\begin{axiom}[GA-4: Effort-Wage Reciprocity]
\label{ax:ga4}
Worker effort responds to wage relative to fair wage:
\begin{equation}
e = e_0 + \gamma \cdot (w - w_{fair})^+ - \lambda \cdot (w_{fair} - w)^+
\end{equation}
where:
\begin{itemize}
    \item $\gamma > 0$: positive reciprocity (gift for gift)
    \item $\lambda > \gamma$: negative reciprocity is stronger (loss aversion)
    \item $(x)^+ = \max(0, x)$
\end{itemize}
\textbf{EBF Connection:} Links to CORE-WHAT fairness dimension and LIT-FEHR reciprocity.
\end{axiom}

\subsection{The Fair Wage-Effort Hypothesis}

Akerlof \& Yellen (1990) formalized how fairness affects unemployment:

\begin{tcolorbox}[colback=green!5!white, colframe=green!75!black]
\textbf{Key Insight:} If workers withdraw effort when paid below fair wage:
\begin{itemize}
    \item Firms prefer to pay $w_{fair}$ rather than face reduced productivity
    \item At $w_{fair} > w_{market-clearing}$, there is involuntary unemployment
    \item This explains wage rigidity and unemployment persistence
\end{itemize}
\end{tcolorbox}

% ============================================================================
% NEAR-RATIONALITY
% ============================================================================
\section{Near-Rationality and Behavioral Macro}

\subsection{Small Deviations, Large Effects}

Akerlof \& Yellen (1985) showed that near-rational behavior can have first-order aggregate effects:

\begin{axiom}[GA-5: Near-Rationality]
\label{ax:ga5}
If agents make small optimization errors $\epsilon$:
\begin{equation}
\text{Individual loss} \propto \epsilon^2 \quad \text{(second-order)}
\end{equation}
\begin{equation}
\text{Aggregate effect} \propto \epsilon \quad \text{(first-order)}
\end{equation}
Small individual irrationalities can cause large macroeconomic fluctuations.
\textbf{EBF Connection:} Justifies behavioral parameters even when ``small.''
\end{axiom}

\subsection{Animal Spirits}

Akerlof \& Shiller (2009) identified five psychological forces driving economies:

\begin{enumerate}
    \item \textbf{Confidence}: Trust and optimism drive investment
    \item \textbf{Fairness}: Perceptions of fair treatment affect wage-setting
    \item \textbf{Corruption}: Bad faith undermines markets
    \item \textbf{Money Illusion}: Nominal vs. real confusion
    \item \textbf{Stories}: Narratives drive economic decisions
\end{enumerate}

\begin{axiom}[GA-6: Confidence Multiplier]
\label{ax:ga6}
Aggregate demand depends on confidence ($C$):
\begin{equation}
Y = \mu(C) \cdot Y_{fundamental}
\end{equation}
where $\mu(C) > 1$ in booms and $\mu(C) < 1$ in busts.
\textbf{EBF Connection:} Provides behavioral foundation for DOMAIN-MACRO.
\end{axiom}

% ============================================================================
% SOCIAL NORMS
% ============================================================================
\section{Social Norms and Custom}

\subsection{Theory of Social Custom}

Akerlof (1980) modeled how social norms affect economic behavior:

\begin{tcolorbox}[colback=yellow!10!white, colframe=yellow!75!black]
\textbf{Social Custom Model:}
\begin{itemize}
    \item Agents gain utility from following norms ($U_{norm}$)
    \item Agents lose utility from being caught violating norms ($U_{reputation}$)
    \item Multiple equilibria possible: high-compliance and low-compliance
    \item Norm changes can be sudden (tipping points)
\end{itemize}
\end{tcolorbox}

\subsection{Connection to Enke's Work}

Akerlof's social norms framework connects to:
\begin{itemize}
    \item \textbf{Enke's Moral Universalism}: Norms vary by kinship structure
    \item \textbf{Tight vs. Loose Cultures}: Norm enforcement varies by society
    \item \textbf{CONTEXT-CULTURE}: Ψ includes norm parameters
\end{itemize}

% ============================================================================
% EBF INTEGRATION
% ============================================================================
\section{EBF Integration}

\subsection{10C Mapping}

\begin{table}[h]
\centering
\begin{tabular}{lll}
\textbf{10C Element} & \textbf{Akerlof Contribution} & \textbf{Parameter} \\
\hline
CORE-WHO (AAA) & Social categories & $c_j$ identity category \\
CORE-WHAT (C) & Identity utility & $I_j$, fairness \\
CORE-HOW (B) & Identity-action complementarity & $\gamma_{identity}$ \\
CORE-WHEN (V) & Procrastination, salience & Time-varying attention \\
CORE-WHERE (BBB) & Information asymmetry α & Lemons parameters \\
CORE-AWARE (AU) & Information completeness & $\alpha$ parameter \\
CORE-READY (AV) & Identity-consistent action & Prescription matching \\
CORE-STAGE (AW) & Identity formation & Category assignment \\
\end{tabular}
\caption{Akerlof contributions mapped to 10C framework}
\end{table}

\subsection{Key Parameter Values for BBB}

\begin{tcolorbox}[colback=green!5!white, colframe=green!75!black, title=Akerlof Parameters for BBB Repository]
\begin{itemize}
    \item $\delta_{identity}$: 0.2-0.5 - Identity loss coefficient
    \item $\gamma_{reciprocity}$: 0.3-0.6 - Positive gift exchange
    \item $\lambda_{negative}$: 0.5-1.0 - Negative reciprocity (stronger)
    \item $w_{premium}$: 10-20\% - Efficiency wage over market
    \item $\epsilon_{near-rational}$: < 0.01 - Near-rationality threshold
    \item $\mu_{confidence}$: 0.8-1.2 - Confidence multiplier range
\end{itemize}
\end{tcolorbox}

% ============================================================================
% WORKED EXAMPLE
% ============================================================================
\section{Worked Example: Identity and Workplace Behavior}

Consider a male nurse (social category conflict):

\textbf{Identity Framework Analysis:}

\begin{align}
U &= U_{material}(wage) + I(actions, category, prescriptions) \\
I &= I_{nurse}(caring, competence) + I_{male}(strength, independence) \\
\text{Conflict}: & \quad P_{nurse} \neq P_{traditional\_male}
\end{align}

\textbf{Prediction:}
\begin{itemize}
    \item If $\delta$ high: Male nurses experience identity tension
    \item May compensate with masculine behaviors (specializing in ER, avoiding pediatrics)
    \item Or: Identity may shift to ``nurse'' as primary category
\end{itemize}

\textbf{Policy Implication:} Changing prescriptions (``men can be nurturing'') reduces identity cost and increases male nursing participation.

% ============================================================================
% SUMMARY
% ============================================================================


%%%%%%%%%%%%%%%%%%%%%%%%%%%%%%%%%%%%%%%%%%%%%%%%%%%%%%%%%%%%%%%%%%%%%%%%%%%%%%%
\section{Key Results}
\label{sec:ga-results}
%%%%%%%%%%%%%%%%%%%%%%%%%%%%%%%%%%%%%%%%%%%%%%%%%%%%%%%%%%%%%%%%%%%%%%%%%%%%%%%

The analysis yields the following key results:

\begin{enumerate}
    \item \textbf{Result 1:} Primary finding with quantitative support
    \item \textbf{Result 2:} Secondary finding with implications
    \item \textbf{Result 3:} Practical application or boundary condition
\end{enumerate}

\section{Summary}

\begin{tcolorbox}[colback=gray!10!white, colframe=gray!75!black, title=Key Takeaways]
\begin{enumerate}
    \item \textbf{Lemons Problem:} Information asymmetry can cause market failure
    \item \textbf{Identity Economics:} Social categories create utility beyond material payoffs
    \item \textbf{Gift Exchange:} Labor relations involve reciprocity, not just market exchange
    \item \textbf{Near-Rationality:} Small behavioral deviations have large macro effects
    \item \textbf{Animal Spirits:} Confidence and narratives drive economic cycles
    \item \textbf{EBF Integration:} Akerlof provides identity utility dimension and fairness foundations
\end{enumerate}
\end{tcolorbox}

% ============================================================================
% GLOSSARY
% ============================================================================
\section{Glossary}

See Appendix G (Glossary) for standard EBF terms. Akerlof-specific terms:

\begin{description}
    \item[Adverse Selection] Market failure when informed party selects against uninformed
    \item[Animal Spirits] Psychological forces (confidence, fairness, stories) driving macro
    \item[Efficiency Wage] Above-market wage paid to elicit effort/loyalty
    \item[Gift Exchange] Reciprocal above-minimum contributions in labor relations
    \item[Identity Utility] Utility derived from matching behavior to social category prescriptions
    \item[Near-Rationality] Small optimization errors with first-order aggregate effects
\end{description}

% ============================================================================
% CRITICAL FOUNDATIONS
% ============================================================================
\section{Critical Foundations}

\begin{enumerate}
    \item Akerlof, G. (1970). The Market for Lemons. \textit{QJE}. [Nobel Prize 2001]
    \item Akerlof, G. \& Kranton, R. (2000). Economics and Identity. \textit{QJE}. [Identity]
    \item Akerlof, G. (1982). Labor Contracts as Partial Gift Exchange. \textit{QJE}. [Gift Exchange]
    \item Akerlof, G. \& Yellen, J. (1990). The Fair Wage-Effort Hypothesis. \textit{QJE}. [Fairness]
    \item Akerlof, G. \& Shiller, R. (2009). Animal Spirits. Princeton UP. [Behavioral Macro]
\end{enumerate}

% ============================================================================
% OPEN ISSUES
% ============================================================================
\section{Open Issues}

\begin{enumerate}
    \item \textbf{Identity measurement:} How to directly measure identity utility in field settings?
    \item \textbf{Category fluidity:} When and how do individuals switch social categories?
    \item \textbf{Prescription formation:} How do category prescriptions emerge and change?
    \item \textbf{Macro foundations:} How to micro-found animal spirits in DSGE models?
\end{enumerate}

% ============================================================================
% REFERENCES
% ============================================================================
\section{References}

\nocite{bcm_master}

For complete bibliography, see Master Bibliography (\texttt{bibliography/master\_references.bib}).

Key Akerlof references integrated in this appendix:
\begin{itemize}[nosep]
    \item akerlof1970lemons, akerlof2000economics, akerlof2010identity
    \item akerlof1982labor, akerlof1990fair, akerlof2009animal
    \item akerlof1985near, akerlof1991procrastination, akerlof1980theory
    \item akerlof2001behavioral\_macro, akerlof2015phishing
\end{itemize}

\end{document}
